\documentclass[11pt,a4paper]{article}
\usepackage[utf8]{inputenc}
\usepackage[T1]{fontenc}
\usepackage[french]{babel}
\usepackage{amsmath}
\usepackage{askme}

% Corrigé — fichier jamais distribué aux élèves.
% Le CLI sync lit les environnements askmecorrection pour remplir la base :
%   python cli/manage.py sync feuille.tex corrige.tex

\begin{document}

\begin{center}
  {\Large \textsc{Corrigé — PGCD et fractions}}
\end{center}

\vspace{10pt}

\begin{askmecorrection}{ex-pgcd-01}
1) On décompose en facteurs premiers :
$120 = 2^3 \times 3 \times 5$ et $84 = 2^2 \times 3 \times 7$.

2) Le PGCD est le produit des facteurs premiers communs, affectés du plus
petit exposant : $\mathrm{PGCD}(120, 84) = 2^2 \times 3 = 12$.

3) On simplifie la fraction par le PGCD :
$\dfrac{84}{120} = \dfrac{84 \div 12}{120 \div 12} = \dfrac{7}{10}$.
\end{askmecorrection}

\begin{askmecorrection}{ex-thales-01}
1) Les droites $(BC)$ et $(DE)$ sont parallèles et les points $A$, $B$, $D$
d'une part, $A$, $C$, $E$ d'autre part sont alignés dans le même ordre.
D'après le théorème de Thalès : $\dfrac{AB}{AD} = \dfrac{AC}{AE}$.

2) On remplace : $\dfrac{4}{6} = \dfrac{6}{AE}$.

3) Donc $AE = \dfrac{6 \times 6}{4} = 9$. Le segment $[AE]$ mesure 9~cm.
\end{askmecorrection}

\end{document}
